%! Absolute Value Equations \& Inequalities — Unit 2, Lesson 2.2 — Homework (Answer Key)
\documentclass[10pt]{article}
\usepackage{algebra2-article}
\usepackage{algebra2-key}   % pulls in -boxes; do NOT also load -boxes
\newcommand{\TallMath}[1]{$\displaystyle #1\rule[-1.4em]{0pt}{3.2em}$}
\newcommand{\numline}[2]{%
  \draw[->, semithick, charcoal] (#1-0.4,0)--(#2+0.4,0);
  \foreach \x in {#1,...,#2}{\draw[charcoal] (\x,0.10)--(\x,-0.10) node[below, font=\scriptsize]{$\x$};}}

\begin{document}

\pageheader{Unit 2, Lesson 2.2}{Homework --- Answer Key}
\namedateperiod

\vspace{0.10in}

\begin{notesbox}{Practice}
Show your work. Isolate the absolute value \emph{before} splitting, and write every inequality
solution in set, interval, and number-line form.

\begin{enumerate}[leftmargin=1.9em, itemsep=11pt]
  \item \textbf{Solve (already isolated).}
    \begin{enumerate}[label=(\alph*), leftmargin=1.8em, itemsep=4pt]
      \item $|x| = 10$: \quad $x = {\color{keyred}\mathbf{10}} \text{ or } {\color{keyred}\mathbf{-10}}$
      \item $|x - 5| = 2$: \quad $x = {\color{keyred}\mathbf{7}} \text{ or } {\color{keyred}\mathbf{3}}$
    \end{enumerate}

  \item \textbf{Isolate, then solve.}
    \begin{enumerate}[label=(\alph*), leftmargin=1.8em, itemsep=4pt]
      \item $|4x + 2| = 10$: \quad $x = {\color{keyred}\mathbf{2}} \text{ or } {\color{keyred}\mathbf{-3}}$
      \item $2|x - 4| - 1 = 9$: \quad $x = {\color{keyred}\mathbf{9}} \text{ or } {\color{keyred}\mathbf{-1}}$
    \end{enumerate}

  \item \textbf{How many solutions?} Write \emph{none}, \emph{one}, or \emph{two}, and give any
        solution(s).
    \begin{enumerate}[label=(\alph*), leftmargin=1.8em, itemsep=4pt]
      \item $|x + 7| = -2$: \ans{none --- a distance can't be negative}
      \item $|x - 3| = 0$: \ans{one, $x = 3$}
      \item $|x| = 5$: \ans{two, $x = 5$ or $-5$}
    \end{enumerate}

  \item \textbf{Solve the ``and'' inequality} $|x + 2| < 6$.
        \\[4pt] Solve: \ans{$-8 < x < 4$} \qquad Set: \ans{$\{x \mid -8 < x < 4\}$} \qquad Interval: \ans{$(-8,\,4)$}
        \begin{center}
        \begin{tikzpicture}[scale=0.5]
          \numline{-9}{5}
          \draw[very thick, forest] (-8,0)--(4,0);
          \draw[very thick, forest, fill=white] (-8,0) circle (3.5pt);
          \draw[very thick, forest, fill=white] (4,0) circle (3.5pt);
        \end{tikzpicture}
        \end{center}

  \item \textbf{Solve the ``or'' inequality} $|2x - 1| \ge 7$.
        \\[4pt] Solve: \ans{$x \le -3$ or $x \ge 4$} \qquad Set: \ans{$\{x \mid x \le -3 \text{ or } x \ge 4\}$} \qquad Interval: \ans{$(-\infty,-3]\cup[4,\infty)$}
        \begin{center}
        \begin{tikzpicture}[scale=0.5]
          \numline{-5}{6}
          \draw[very thick, forest, <-] (-5.4,0)--(-3,0);
          \draw[very thick, forest, ->] (4,0)--(6.4,0);
          \fill[forest] (-3,0) circle (3.5pt); \fill[forest] (4,0) circle (3.5pt);
        \end{tikzpicture}
        \end{center}

  \item \textbf{Model it.} A thermostat holds a room within $2^\circ$F of the $68^\circ$F setting.
    \begin{enumerate}[label=(\alph*), leftmargin=1.8em, itemsep=4pt]
      \item Write an absolute-value inequality for an acceptable temperature $T$. \ans{$|T - 68| \le 2$}
      \item Solve it and state the coolest and warmest acceptable temperatures. \ans{$66 \le T \le 70$; coolest $66^\circ$F, warmest $70^\circ$F}
    \end{enumerate}

  \item \textbf{Explain.} In your own words, how do you decide whether an absolute-value inequality
        becomes an ``and'' (one interval) or an ``or'' (two rays)?
        \ansline{``$<$/$\le$'' means \emph{within} a distance $\to$ ``and'' (one interval between the bounds);}
        \ansline{``$>$/$\ge$'' means \emph{beyond} a distance $\to$ ``or'' (two rays outside the bounds).}
\end{enumerate}
\end{notesbox}

\begin{extensionbox}
\textbf{Tolerance and reasoning.}
\begin{enumerate}[label=(\alph*), leftmargin=1.9em, itemsep=6pt]
  \item A bolt is acceptable if its diameter $d$ is within $0.2$ mm of the $15$ mm target. Write and
        solve an inequality; give the smallest and largest acceptable diameters.
  \item For what values of $k$ does $|x - 3| < k$ have \emph{no solution}? For what values of $k$ is
        \emph{every} real number a solution of $|x - 3| > k$? Justify from the meaning of absolute
        value.
\end{enumerate}
\ansline{(a) $|d - 15| \le 0.2 \Rightarrow 14.8 \le d \le 15.2$; smallest $14.8$ mm, largest $15.2$ mm.}
\ansline{(b) $|x-3|<k$ has no solution when $k \le 0$; $|x-3|>k$ holds for all reals when $k < 0$.}
\ansline{A distance is always $\ge 0$: it can't be less than $0$ or a negative, and it always exceeds a negative.}
\end{extensionbox}

\begin{spiralbox}
\textbf{Coming next --- Lesson 2.3: Absolute Value Functions \& Transformations.} Today you
\emph{solved} $|ax+b|=c$; next you'll \emph{graph} $y=|ax+b|$ --- a \textbf{V-shape}. The two
solutions you found for an equation are exactly where that V crosses the horizontal line $y=c$, which
previews the vertex, axis of symmetry, and minimum you'll study formally in Unit~3.
\end{spiralbox}

\begin{teachernote}
Watch item~2: students must isolate ($|4x+2|=10$, $|x-4|=5$) before splitting --- a very common miss.
Item~3 is a quick check that ``distance can't be negative'' (none), ``distance zero'' (one), and the
generic case (two) are all understood conceptually. In item~6, ``within'' includes the endpoints, so
it is ``$\le$''; students who write ``$<$'' exclude $66^\circ$ and $70^\circ$. Extension (b)'s boundary
$k=0$ is worth a mention: $|x-3|>0$ excludes $x=3$, so ``all reals'' needs $k<0$ strictly.
\end{teachernote}

\end{document}
